\section{方法的可复用部分与下一步问题}\label{sec:methods}

\subsection{从定位损失到选择方法}
研究的起点是把一个全范围问题拆成可测量的若干损失：零点贡献、截断误差、非素数素数幂以及端点贡献。若归一化后这些量的严格上界之和小于 $1$，主项就留下正的素数贡献。这个框架的作用在于辨认困难来自哪里；提高小数精度只能减少数值包围误差，不能缩小一个本来很大的解析上界。

早期围绕 $k=4$ 的试验揭示了中段与尾段的区别。尾段可以利用随 $L$ 衰减的指数因子建立解析单调性；中段则对密度常数、权函数的变换衰减和有限高度的拆分更加敏感。因而本项目把 $k=6$ 作为能够闭合全部范围的阶段成果，同时保留对较小指数的局部计算和结构分析。

梯形权的选择兼顾两个要求：权非负且在严格开区间之外为零，使正加权和能提取素数；其四点差分表达又允许从同一个积分显式公式得到零次、一次和二次三种高度衰减界。取三者的最小值，可以在不同高度利用合适的估计。更高平滑阶数并不自动更好，因为主项质量、导数范数和相对区间长度必须同时计算。

\subsection{形式化促成的数学重组}
原始书面证明逐个认证中段和远段单元；远段还采用另一种多项式权。最终形式化把主要有限范围重组为两个基准证书、一个连续传播定理和一个有限高度 RH 比较。这样既减少重复的超越估计，也给出了对所有实参数成立的结构性解释。

连续传播的关键不是把一个成功参数复制到邻近采样点，而是同时平移 $A$ 和 $u$、缩放 $q$。核的高度质量保持不变，密度包围则由导数上界控制为不增加。第\ref{sec:coverage}节证明的适用条件包含 $k\ge4$；这说明该引理本身能服务于多个指数，但全范围素数结论还要求某个合法起点具有严格小于 $1$ 的完整预算，并与整数下段衔接。

有限高度 RH 的使用也有明确目的。在 $L$ 向下移动时，一般零点密度预算会变坏；低于 $H=3\cdot10^{12}$ 的零点却可以使用 $\beta=1/2$ 的精确信息。将低高度单独扣出后，高高度核的缩放补偿了部分损失，留下可由三档有理系数认证的比较。

\subsection{计算余量应如何解释}
预算小于 $1$ 的差额是所用估计的确定性余量。例如下中段的保守上界为 $0.966283$，因而这一预算还留有 $0.033717$。这表示在当前分项和证明条件下，扣除项严格不足主项；它不是命题正确的概率，也不允许把每个上游参数任意改动同一百分比。

原始 $94$ 单元证书另验证过全部三列 CS 系数同时增加 $6\%$ 的敏感性试验。该结果属于原始证书与其固定包围。新下中段采用另一套保守比较，故应分别报告其实际预算，不能把原敏感性结论直接迁移给当前 Lean 定理。研究历史保留两套记录及其明确对应的数学版本。

CHL 的舍入参数复现问题同样要区分目标常数与见证参数。我们在 H7 中接受的是指定短区间结论；作者表中有限位的辅助参数是建立它的见证。六行印刷参数生成的 $\Delta$ 略低于报告目标，微调参数后的普通数值试验显示仍有可行余量。严格区间诊断已经排除了浮点伪象，但完整判据的独立认证仍是一项清晰、有限的后续工作。这项问题保留在输入登记中，而不改变本次给定 H1--H7 的逻辑推导。

\subsection{通向较小指数的任务}
若考虑一般长度 $h=kx^{1-1/k}$，则控制核的参数具有 $q\asymp x^{1/k}/k$ 的尺度。减小 $k$ 会增大相对所需截断高度，并使同一 $L$ 下的加权零点预算更难压低。可复用的部分包括权函数恒等式、密度分层、高度尾项、整数矩认证以及部分连续传播引理；需要重新求解的部分包括基准参数、全区间预算、下端覆盖与两者之间的接合。

因此，对 $k=4$，下一步应寻找能够与已知整数下段相接的有效基准区间，并定位阻碍预算下降的密度区间与高度区间。对 $k=3$，现有连续传播定理的 $k\ge4$ 条件本身还需改进。局部成功和一般形式的引理都是可积累成果，只有所有范围闭合后才升级为新的全范围结论。

本文的已证明成果、历史探索及未来任务在仓库中分层保存。这样的记录允许后续研究复用现有组件，而不把某次未闭合的探索误读为主定理的一部分。
